\section{Prompt Engineering：提示工程}

\subsection{Role Prompting（角色提示）}

Role Prompting 是最简单但极其有效的 prompt 技巧。讲师展示了一个精彩的例子：

\begin{quote}
\textbf{Prompt 1}：What is $100 \times 100 \div 400 \times 56$? \\
\textbf{模型回答}：280 \quad {\color{red}（错误）}
\end{quote}

\begin{quote}
\textbf{Prompt 2}：You are an MIT mathematician. What is $100 \times 100 \div 400 \times 56$? \\
\textbf{模型回答}：1400 \quad {\color{green!60!black}（正确）}
\end{quote}

仅仅添加了 ``You are an MIT mathematician'' 这几个词，模型就从给出错误答案变为给出正确答案。

讲师对此的直觉解释是：

\begin{importantbox}{Role Prompting 为何有效——概率空间的条件化}
模型的训练数据来自互联网。互联网上大量用户在回答数学问题时给出了错误答案。但如果你限定到那些以``我是 MIT 数学家''开头回答问题的人群，他们给出正确答案的概率就大幅提高。Role Prompting 本质上是对模型的输出概率分布做了条件化（conditioning），将概率质量集中到训练数据中高质量回答所在的区域。
\end{importantbox}

讲师幽默地补充道：由于 word embeddings 的模糊性，``MIT mathematician'' 可能也激活了 ``Harvard mathematician'' 所在的嵌入空间区域——因此这个技巧对 Harvard 的数学家也有效。

\subsection{Chain-of-Thought Prompting（思维链提示）}

Chain-of-Thought（CoT）Prompting 是一种引导模型``展示解题步骤''的技巧。

\textbf{标准 Prompting}：给一个数学应用题和答案作为 one-shot 示例，然后给一个新问题。模型直接输出（错误的）答案。

\textbf{CoT Prompting}：同样给一个数学应用题，但示例中包含完整的推理步骤，然后给新问题。模型会模仿这种``展示步骤''的模式，在给出最终答案前逐步推导，最终更可能得到正确答案。

\begin{importantbox}{``Let's think step by step'' 的魔力}
后续研究表明，一个令人惊讶的发现是：你甚至不需要给出完整的推理示例，只需在 prompt 末尾添加一句 ``Let's think step by step''，就足以触发模型的逐步推理模式，大幅提升数学和逻辑任务的准确率。
\end{importantbox}

讲师对 CoT 有效性的直觉解释基于\textbf{错误驱动学习}（error-driven learning）的视角：

\begin{knowledgebox}{CoT 有效性的直觉解释——错误表面积}
在标准 prompting 下，一个数学应用题的输出可能是 ``The answer is 27''。在整个输出中，模型可能出错的地方只有最后的数字 ``27''——错误表面积（error surface area）很小，模型在训练时通过反向传播获得的学习信号也很有限。

当使用 CoT 时，模型需要生成完整的中间步骤，如 ``There are 16 players... half of them quit, so 8 remain... a quarter of those lost, so 2 lost... 8 - 2 = 6''。每个中间步骤都是一个潜在的出错点。在训练阶段，这种更长的输出序列提供了更多的错误表面积，模型有更多机会犯错并据此更新权重，从而学会更准确的推理。

注意：这个学习过程发生在\textbf{训练阶段}，推理时模型的权重并不更新。CoT 之所以在推理时有效，是因为模型在训练时已经从这类逐步推理的数据中学到了相应的能力。
\end{knowledgebox}

\subsection{本章小结}

Prompt Engineering 的两大核心技巧——Role Prompting 和 Chain-of-Thought Prompting——本质上都是在操纵模型的概率分布。Role Prompting 通过身份设定条件化输出分布；CoT 通过要求逐步推理增加模型的``思考''空间。两者可以结合使用以获得最佳效果。

